%% kapa: exam=1 round=37 year=2026 period=1 form=A subject=civil_law subject_ko=민법 seq=37 number=37
%% source: Q-Net 감정평가사 기출문제 (cst00309) artlSeq=5258924 file=raw/1008/5258924/2026년 제37회 감정평가사 1차 1교시.pdf
%% conversion: kordoc 4.24.51 -> tools/gen1.py (auto); review=proofread-sonnet; flags=none
\begin{question}{37}
민사유치권에 관한 설명으로 옳은 것을 모두 고른 것은? (다툼이 있으면 판례에 따름)
\begin{framedtext}
ㄱ. 채권과 물건 사이에 견련관계가 인정되는 경우에는 그 채무불이행으로 인한 손해배상채권과 물건 사이에도 견련관계가 인정될 수 있다. \par
ㄴ. 정당한 이유가 있는 때에는 유치권자가 유치물로 직접 변제에 충당할 것을 법원에 청구할 수 있으며, 이 경우 유치권자는 미리 채무자에게 통지하지 않아도 된다. \par
ㄷ. 점유를 침탈당한 유치권자가 점유회수의 소를 제기하여 승소판결을 받았더라도 점유를 회복하지 않으면 유치권이 부활하지 않는다. \par
ㄹ. 유치권자의 선관주의의무 위반을 이유로 한 유치권소멸청구권은 특별한 사정이 없는 한 유치권자의 선관주의의무 위반이 있은 이후에 유치물의 소유권을 취득한 제3자도 이를 행사할 수 있다. \par
\end{framedtext}
\begin{choices}
  \item ㄱ, ㄴ
  \item ㄱ, ㄷ
  \item ㄱ, ㄴ, ㄷ
  \item ㄱ, ㄷ, ㄹ
  \item ㄴ, ㄷ, ㄹ
\end{choices}
\end{question}
\answer{4}
